% ════════════════════════════════════════════════════════════════════════════
%  preamble.tex
%  Преамбула: пакеты, тема, цвета, стили tcolorbox, шапка/подвал.
%  Менять только если нужно изменить ОФОРМЛЕНИЕ. Контент презентации
%  лежит в main.tex (метаданные) и в slides/ (сами слайды).
% ════════════════════════════════════════════════════════════════════════════

% ─── Пакеты ─────────────────────────────────────────────────────────────────
\usepackage[utf8]{inputenc}
\usepackage[T2A]{fontenc}
\usepackage[russian]{babel}
\usepackage{booktabs}
\usepackage{array}
\usepackage{colortbl}
\usepackage{pgfplots}
\pgfplotsset{compat=1.18}
\usepackage{tikz}
\usetikzlibrary{shapes.geometric, arrows.meta, positioning, shadows, fit, calc}
\usepackage{amsmath}
\usepackage{amssymb}
\usepackage{fontawesome5}
\usepackage{multirow}
\usepackage{makecell}
\usepackage{xcolor}
\usepackage{tcolorbox}
\tcbuselibrary{skins, breakable}

% ─── Тема и цветовая схема ──────────────────────────────────────────────────
% TODO (опционально): чтобы поменять «фирменный» цвет всей презентации,
% достаточно изменить RGB у primary / secondary / accent ниже.
\usetheme{Madrid}
\usecolortheme{default}

\definecolor{primary}{RGB}{13, 71, 161}     % основной (синий)
\definecolor{secondary}{RGB}{21, 101, 192}  % вторичный
\definecolor{accent}{RGB}{255, 143, 0}      % акцент (оранжевый)
\definecolor{success}{RGB}{27, 94, 32}      % зелёный (успех)
\definecolor{danger}{RGB}{183, 28, 28}      % красный (предупреждение)
\definecolor{light}{RGB}{227, 242, 253}     % светлый фон блоков
\definecolor{darktext}{RGB}{13, 13, 13}
\definecolor{mygray}{RGB}{96, 125, 139}

\setbeamercolor{palette primary}{bg=primary, fg=white}
\setbeamercolor{palette secondary}{bg=secondary, fg=white}
\setbeamercolor{palette tertiary}{bg=primary!80, fg=white}
\setbeamercolor{palette quaternary}{bg=primary, fg=white}
\setbeamercolor{structure}{fg=primary}
\setbeamercolor{frametitle}{bg=primary, fg=white}
\setbeamercolor{title}{bg=primary, fg=white}
\setbeamercolor{block title}{bg=primary, fg=white}
\setbeamercolor{block body}{bg=light}
\setbeamercolor{block title alerted}{bg=danger, fg=white}
\setbeamercolor{block body alerted}{bg=danger!10}
\setbeamercolor{block title example}{bg=success, fg=white}
\setbeamercolor{block body example}{bg=success!10}
\setbeamercolor{item}{fg=primary}
\setbeamercolor{subitem}{fg=secondary}

\setbeamerfont{title}{size=\Large, series=\bfseries}
\setbeamerfont{frametitle}{size=\large, series=\bfseries}

\setbeamertemplate{navigation symbols}{}

% ─── Подвал слайда ──────────────────────────────────────────────────────────
\setbeamertemplate{footline}{%
  \leavevmode%
  \hbox{%
    \begin{beamercolorbox}[wd=.4\paperwidth,ht=2.5ex,dp=1.125ex,leftskip=.3cm]{palette primary}%
      \usebeamerfont{author in head/foot}\insertshortauthor
    \end{beamercolorbox}%
    \begin{beamercolorbox}[wd=.4\paperwidth,ht=2.5ex,dp=1.125ex,center]{palette secondary}%
      \usebeamerfont{title in head/foot}\insertshorttitle
    \end{beamercolorbox}%
    \begin{beamercolorbox}[wd=.2\paperwidth,ht=2.5ex,dp=1.125ex,rightskip=.3cm,right]{palette primary}%
      \usebeamerfont{date in head/foot}
      \insertframenumber{} / \inserttotalframenumber
    \end{beamercolorbox}%
  }%
}

% ─── Стили цветных «карточек» (tcolorbox) ───────────────────────────────────
% В слайдах используются 4 готовых стиля:
%   infobox    — синяя полоса слева (нейтральная информация)
%   formulabox — оранжевая рамка   (формулы / ключевые определения)
%   warnbox    — красная полоса    (предупреждения, минусы)
%   successbox — зелёная полоса    (плюсы, выводы)
\tcbset{
  infobox/.style={
    enhanced, boxrule=0pt, frame hidden,
    colback=light, left=4pt, right=4pt, top=2pt, bottom=2pt,
    borderline west={3pt}{0pt}{primary}
  },
  formulabox/.style={
    enhanced, boxrule=1pt,
    colframe=accent, colback=accent!10,
    left=4pt, right=4pt, top=2pt, bottom=2pt,
    rounded corners
  },
  warnbox/.style={
    enhanced, boxrule=0pt, frame hidden,
    colback=danger!8, left=4pt, right=4pt, top=2pt, bottom=2pt,
    borderline west={3pt}{0pt}{danger}
  },
  successbox/.style={
    enhanced, boxrule=0pt, frame hidden,
    colback=success!8, left=4pt, right=4pt, top=2pt, bottom=2pt,
    borderline west={3pt}{0pt}{success}
  }
}
